\documentclass[11pt]{article}

\usepackage[a4paper,margin=1in]{geometry}
\usepackage{amsmath,amssymb,mathtools}
\usepackage{booktabs,tabularx,array,longtable}
\usepackage{enumitem}
\usepackage[hidelinks]{hyperref}
\usepackage[T1]{fontenc}

\newcommand{\R}{\mathbb{R}}
\newcommand{\N}{\mathbb{N}}
\newcommand{\eps}{\varepsilon}
\newcommand{\Sym}{\operatorname{Sym}}
\newcommand{\lb}{\operatorname{lb}}
\newcommand{\ub}{\operatorname{ub}}
\newcommand{\midp}{\operatorname{mid}}
\newcommand{\rad}{\operatorname{rad}}
\newcommand{\wid}{\operatorname{wid}}
\newcommand{\magI}{\operatorname{mag}}
\newcommand{\mig}{\operatorname{mig}}
\newcommand{\Mean}{\operatorname{Mean}}
\newcommand{\Max}{\operatorname{Max}}
\newcommand{\NA}{\texttt{NA}}

\title{Diagnostics and Metrics for Polynomial-Zonotope\\
Neural-Network Certification}
\author{}
\date{}

\begin{document}
\maketitle

\begin{abstract}
This note fixes a coherent and explicitly computable set of diagnostics for
certified neural networks in PDE applications.  It is compatible with the
companion reference \emph{Notation and Terminology for Polynomial-Zonotope
Neural-Network Certification}.  In particular, there are exactly two classes
of noise symbols: domain noise symbols \(\alpha\) and approximation noise
symbols \(\eta\), with combined vector \(\eps=(\alpha,\eta)\).  The word
\emph{residual} is reserved for an actual function difference, such as an
activation-approximation residual or a PDE residual, and is not used as an
alternative name for \(\eta\).

The note gives mathematical definitions, aggregation rules, canonical output
schemas, and two required benchmark levels.  The \emph{mini benchmark} is the
smallest standardized report suitable for routine method comparisons.  The
\emph{medium benchmark} contains the complete mini benchmark and additionally
records per-neuron preactivation intervals, activation approximation-error
radii, and layerwise normalized interval-hull radii.
\end{abstract}

\section{Compatibility with the notation reference}

Let
\[
 \Phi:\mathcal X\subseteq\R^{d_0}\longrightarrow\R^{d_{\rm out}}
\]
be the complete neural network and let
\[
 X(\alpha),\qquad \alpha\in[-1,1]^p,
\]
be a polynomial-zonotope parametrization of the physical input set
\(\mathcal X\).  The network is written as
\[
 \Phi=A_L\circ\sigma_L\circ\cdots\circ A_1\circ\sigma_1\circ A_0,
 \qquad A_\ell(y)=W_\ell y+b_\ell.
\]
At hidden layer \(\ell\), \(z_\ell\) and \(y_\ell\) denote the preactivation
and postactivation, while \(Z_\ell(\eps)\) and \(Y_\ell(\eps)\) denote their
polynomial enclosures.  The corresponding polynomial enclosures of the first
and second physical-input derivatives are denoted by \(J_\ell(\eps)\) and
\(H_\ell(\eps)\).

For the complete network, write
\[
 Y(\eps),\qquad J(\eps),\qquad H(\eps)
\]
for the final polynomial enclosures of
\[
 \Phi(X(\alpha)),\qquad
 \left.\partial\Phi\right|_{X(\alpha)},\qquad
 \left.\partial^2\Phi\right|_{X(\alpha)}.
\]
The derivatives are always with respect to the physical input \(x\).  The
variables \(\alpha\) and \(\eta\) only parametrize the certified enclosure.

For \(r\in\{0,1,2\}\), with \(\sigma^{(0)}=\sigma\), the canonical scalar
activation enclosure is
\[
 \sigma^{(r)}(t)
 \in \widehat\sigma_{\ell i}^{(r)}(t)
      +\rho_{\ell i}^{(r)}\eta_{\ell i}^{(r)},
 \qquad \eta_{\ell i}^{(r)}\in[-1,1].
\]
Here \(\rho_{\ell i}^{(r)}\geq0\) is the
\emph{approximation-error radius}, \(\eta_{\ell i}^{(r)}\) is an
\emph{approximation noise symbol}, and
\[
 r_{\ell i}^{(r)}(t)
 =\sigma^{(r)}(t)-\widehat\sigma_{\ell i}^{(r)}(t)
\]
is the actual residual function.

\section{Cells, indices, and physical-volume weights}

Let
\[
 \mathcal X=\bigcup_{c=1}^{N_{\rm cell}}\mathcal X_c
\]
be the partition used for certification, with pairwise disjoint interiors.
For cell \(c\), let \(X_c(\alpha)\) denote its domain parametrization.  We use
\[
 c=1,\ldots,N_{\rm cell},\qquad
 i=1,\ldots,d_{\rm out},\qquad
 a,b=1,\ldots,d_0.
\]
Thus \(i\) indexes the output component, \(a,b\) index physical-input
coordinates, and \(c\) indexes the certification cell.

For unequal physical cell volumes, define
\[
 \omega_c
 =\frac{|\mathcal X_c|}{\sum_{k=1}^{N_{\rm cell}}|\mathcal X_k|},
 \qquad \sum_c\omega_c=1.
\]
When all cells have equal volume, \(\omega_c=1/N_{\rm cell}\).  If physical
cell volumes are unavailable, the resulting average must be labelled an
\emph{unweighted cell average}; it must not be called a domain average.

\section{Interval terminology}

Let \(I=[\underline I,\overline I]\) be a nonempty closed real interval.
Define
\[
 \lb(I)=\underline I,\qquad \ub(I)=\overline I,
\]
\[
 \midp(I)=\frac{\underline I+\overline I}{2},\qquad
 \rad(I)=\frac{\overline I-\underline I}{2},
\]
\[
 \wid(I)=\overline I-\underline I=2\rad(I),
\]
\[
 \magI(I)=\max\{|\underline I|,|\overline I|\},
\]
and
\[
 \mig(I)=
 \begin{cases}
  \min\{|\underline I|,|\overline I|\},&0\notin I,\\
  0,&0\in I.
 \end{cases}
\]
The magnitude is a certified upper bound on \(|z|\) for every \(z\in I\),
while the mignitude is the largest certified lower bound on \(|z|\) over the
interval.

\subsection{Final interval hulls}

For every cell \(c\), intervalize the final polynomial enclosures over all
domain and approximation noise symbols:
\[
 [Y^i]_c=[\underline Y_c^i,\overline Y_c^i],
\]
\[
 [J_{ia}]_c=[\underline J_{c,ia},\overline J_{c,ia}],
\]
\[
 [H_{iab}]_c=[\underline H_{c,iab},\overline H_{c,iab}].
\]
These intervals enclose, respectively,
\[
 \Phi^i(X_c(\alpha)),\qquad
 \left.\partial_a\Phi^i\right|_{X_c(\alpha)},\qquad
 \left.\partial_{ab}\Phi^i\right|_{X_c(\alpha)}.
\]
Unless explicitly stated otherwise, all final width metrics are computed from
these interval hulls.

\section{Scalar tightness metrics}

For a scalar interval \(I\), define the absolute width and radius by
\[
 w(I)=\wid(I),\qquad r(I)=\rad(I).
\]
The absolute width is the primary dimensional tightness metric.

The local relative radius is
\[
 \rho_{\rm loc}(I)=
 \begin{cases}
  \dfrac{\rad(I)}{\magI(I)},&\magI(I)>0,\\[2mm]
  0,&I=[0,0].
 \end{cases}
\]
No stabilization parameter is used.  This metric must always be reported
alongside an absolute metric because an interval can be very small in absolute
terms while still having local relative radius near one when it is centered
near zero.

The local relative width is
\[
 \rho_{\rm wid}(I)=
 \begin{cases}
  \dfrac{\wid(I)}{\magI(I)},&\magI(I)>0,\\[2mm]
  0,&I=[0,0].
 \end{cases}
\]
It is dimensionless and explicitly computable, and it satisfies
\[
 \rho_{\rm wid}(I)=2\rho_{\rm loc}(I).
\]
Thus local relative width and local relative radius contain the same
information.  The glossary includes both terms for completeness, but the
canonical compact summaries use the radius form to avoid redundant headline
metrics.  For intervals crossing zero, the local relative width may be as large
as two.

The sign-certification indicator is
\[
 s(I)=
 \begin{cases}
  1,&0\notin I,\\
  0,&0\in I.
 \end{cases}
\]

\subsection{Familywise reference scales and normalized radii}

Define the data-derived scales
\[
 S_Y=\max_{c,i}\magI([Y^i]_c),
\]
\[
 S_J=\max_{c,i,a}\magI([J_{ia}]_c),
\]
\[
 S_H=\max_{c,i,a,b}\magI([H_{iab}]_c).
\]
For a family \(Q\in\{Y,J,H\}\) and an interval \(I\) from that family,
define
\[
 \rho_{\rm glob}(I;S_Q)=
 \begin{cases}
  \dfrac{\rad(I)}{S_Q},&S_Q>0,\\[2mm]
  0,&S_Q=0.
 \end{cases}
\]
This is the canonical \emph{global normalized radius}.  Equivalently,
\(\wid(I)/(2S_Q)\) gives the same value.  Only the radius form is used in the
canonical schemas.

\paragraph{Scope of a reference scale.}
A reported value of \(S_Y,S_J,S_H\) is valid only for the exact collection of
cells and components summarized in the same output record.  When comparing
methods, the preferred comparison uses a shared reference scale computed from
a designated reference family, or reports absolute widths in addition to each
method's own normalized radius.

\section{Aggregation rules}

Let \(q_{c,k}\geq0\) be a cellwise diagnostic, where \(k\) collects all
non-cell indices.  Its volume-weighted mean is
\[
 \Mean_\omega(q)
 =\sum_{c=1}^{N_{\rm cell}}\omega_c
   \left(\frac1{N_k}\sum_{k=1}^{N_k}q_{c,k}\right).
\]
Its maximum is
\[
 \Max(q)=\max_{c,k}q_{c,k}.
\]
The canonical empirical quantiles are
\[
 Q_{0.50}(q),\qquad Q_{0.90}(q),\qquad Q_{0.99}(q),
\]
computed after flattening the finite collection \(\{q_{c,k}\}_{c,k}\).  The
implementation must record the interpolation convention used by its numerical
library.  The default recommendation is the library's linear interpolation
convention.

The worst-cell index is
\[
 c_{\max}(q)
 =\min\operatorname*{arg\,max}_c\left(\max_k q_{c,k}\right),
\]
where the minimum resolves ties deterministically.

\section{Function-value, Jacobian, and Hessian diagnostics}

\subsection{Function values}

For every \(c,i\), define
\[
 w^Y_{c,i}=\wid([Y^i]_c),\qquad
 \rho_{c,i}^{Y,{\rm loc}}=\rho_{\rm loc}([Y^i]_c),\qquad
 \rho_{c,i}^{Y,{\rm glob}}=\rho_{\rm glob}([Y^i]_c;S_Y).
\]
The required summaries are the volume-weighted mean, maximum, and
\(0.50,0.90,0.99\) quantiles of \(w^Y_{c,i}\), together with the
volume-weighted mean and maximum of \(\rho_{c,i}^{Y,{\rm glob}}\).

Define the componentwise width vector and its norms by
\[
 w_c^Y=(w^Y_{c,1},\ldots,w^Y_{c,d_{\rm out}}),
\]
\[
 W_{c,2}^Y=\|w_c^Y\|_2,\qquad
 W_{c,\infty}^Y=\|w_c^Y\|_\infty.
\]

\subsection{Jacobians}

For every \(c,i,a\), define
\[
 w^J_{c,ia}=\wid([J_{ia}]_c),\qquad
 \rho_{c,ia}^{J,{\rm loc}}=\rho_{\rm loc}([J_{ia}]_c),\qquad
 \rho_{c,ia}^{J,{\rm glob}}=\rho_{\rm glob}([J_{ia}]_c;S_J).
\]
Let \(W_c^J=(w^J_{c,ia})_{i,a}\).  The Jacobian Frobenius width is
\[
 W_{c,F}^J=\|W_c^J\|_F
 =\left(\sum_{i=1}^{d_{\rm out}}\sum_{a=1}^{d_0}
 (w^J_{c,ia})^2\right)^{1/2}.
\]
For output component \(i\), the gradient width is
\[
 W_{c,i}^{\nabla}
 =\left(\sum_{a=1}^{d_0}(w^J_{c,ia})^2\right)^{1/2}.
\]
For physical-input coordinate \(a\), define
\[
 W_{c,a}^{\rm input}
 =\left(\sum_{i=1}^{d_{\rm out}}(w^J_{c,ia})^2\right)^{1/2}.
\]
The required summaries are the same entrywise quantities as for function
values.  In addition, report
\[\Mean_\omega(W_F^J)\qquad\text{and}\qquad \max_c W_{c,F}^J.\]

\subsection{Hessians}

For every \(c,i,a,b\), define
\[
 w^H_{c,iab}=\wid([H_{iab}]_c),\qquad
 \rho_{c,iab}^{H,{\rm loc}}=\rho_{\rm loc}([H_{iab}]_c),\qquad
 \rho_{c,iab}^{H,{\rm glob}}=\rho_{\rm glob}([H_{iab}]_c;S_H).
\]
The Hessian Frobenius width is
\[
 W_{c,F}^H
 =\left(\sum_{i=1}^{d_{\rm out}}\sum_{a=1}^{d_0}
 \sum_{b=1}^{d_0}(w^H_{c,iab})^2\right)^{1/2}.
\]
If only entries with \(a\leq b\) are stored, the implementation must either
reconstruct the full symmetric tensor or use weight one for \(a=b\) and weight
two for \(a<b\).

The Hessian symmetry discrepancy is
\[
 E_{\rm sym}^H
 =\max_{c,i,a,b}\max\left\{
 |\underline H_{c,iab}-\underline H_{c,iba}|,
 |\overline H_{c,iab}-\overline H_{c,iba}|
 \right\}.
\]

\section{Certified function-space norm metrics}

Let \(\mathcal N\geq0\) be a norm or squared norm with certified interval
\[
 [\mathcal N]=[\underline{\mathcal N},\overline{\mathcal N}],
 \qquad 0\leq\underline{\mathcal N}\leq\overline{\mathcal N}.
\]
Always report the raw interval and its absolute width
\[
 w_{\mathcal N}=\overline{\mathcal N}-\underline{\mathcal N}.
\]
The canonical relative norm width is
\[
 \rho_{\mathcal N}=
 \begin{cases}
  \dfrac{\overline{\mathcal N}-\underline{\mathcal N}}
        {\overline{\mathcal N}},&\overline{\mathcal N}>0,\\[2mm]
  0,&\overline{\mathcal N}=0.
 \end{cases}
\]
It lies in \([0,1]\).  The equivalent lower-bound informativeness is
\[
 \gamma_{\mathcal N}=
 \begin{cases}
  \dfrac{\underline{\mathcal N}}{\overline{\mathcal N}},
  &\overline{\mathcal N}>0,\\[2mm]
  1,&\overline{\mathcal N}=0,
 \end{cases}
 \qquad \gamma_{\mathcal N}=1-\rho_{\mathcal N}.
\]
Only \(\rho_{\mathcal N}\) is required as the canonical tightness
summary.  For a nonnegative norm interval, this relative norm width is exactly
the local relative width defined above because
\(\magI([\underline{\mathcal N},\overline{\mathcal N}])
=\overline{\mathcal N}\).

\subsection{Domain-volume-normalized norm intervals}

Let
\[
 V_{\mathcal X}=|\mathcal X|>0
\]
be the physical integration-domain volume.  The
\emph{domain-volume-normalized norm} is
\[
 \mathcal N_{\rm vol}=\frac{\mathcal N}{\sqrt{V_{\mathcal X}}}.
\]
For a certified norm interval, define
\[
 [\mathcal N_{\rm vol}]
 =\left[
 \frac{\underline{\mathcal N}}{\sqrt{V_{\mathcal X}}},
 \frac{\overline{\mathcal N}}{\sqrt{V_{\mathcal X}}}
 \right].
\]
For a certified squared-norm interval
\([\mathcal N^2]=[\underline{\mathcal N^2},
\overline{\mathcal N^2}]\), define
\[
 [\mathcal N_{\rm vol}^2]
 =\left[
 \frac{\underline{\mathcal N^2}}{V_{\mathcal X}},
 \frac{\overline{\mathcal N^2}}{V_{\mathcal X}}
 \right].
\]
For the \(L^2\) norm, \(\mathcal N_{\rm vol}\) is the root-mean-square
magnitude of the function over \(\mathcal X\).  For \(W^{1,2}\) and
\(W^{2,2}\), it is the square root of the mean value of the corresponding
value-and-derivative energy density.  The term \emph{RMS norm} is canonical
only for \(L^2\); for Sobolev norms use \emph{domain-volume-normalized
Sobolev norm}.

This metric is not a tightness metric.  It removes the trivial
\(\sqrt{V_{\mathcal X}}\) scaling of an \(L^2\)-type norm and is useful
when comparing experiments on domains with different volumes.  It does not
normalize by the size of the represented function and does not measure
interval overestimation.  Positive rescaling leaves the relative norm width
unchanged:
\[
 \frac{\wid([\mathcal N_{\rm vol}])}
 {\ub([\mathcal N_{\rm vol}])}
 =
 \frac{\wid([\mathcal N])}{\ub([\mathcal N])}
 =\rho_{\mathcal N}.
\]
Consequently, the domain-volume-normalized interval and the relative norm
width must be reported as separate metrics.  The word \emph{normalized}
without a qualifier must not be used for either one.

For \(L^2\)- and Sobolev norms, retain both the squared-norm interval and the
norm interval.  The primary relative-width diagnostic is computed before the
square root.

Define
\[
 I_0=\int_{\mathcal X}|\Phi(x)|_2^2\,dx,
\]
\[
 I_1=\int_{\mathcal X}\left|\left.\partial\Phi\right|_x\right|_F^2\,dx,
\]
\[
 I_2=\int_{\mathcal X}\left|\left.\partial^2\Phi\right|_x\right|_F^2\,dx.
\]
Then
\[
 \|\Phi\|_{L^2}^2=I_0,\qquad
 \|\Phi\|_{W^{1,2}}^2=I_0+I_1,
\]
\[
 \|\Phi\|_{W^{2,2}}^2=I_0+I_1+I_2.
\]
For the terms present in a chosen norm, define
\[
 C_k^{\rm upper}
 =\begin{cases}
 \dfrac{\overline I_k}{\sum_j\overline I_j},
 &\sum_j\overline I_j>0,\\
 0,&\sum_j\overline I_j=0,
 \end{cases}
\]
and
\[
 C_k^{\rm width}
 =\begin{cases}
 \dfrac{\overline I_k-\underline I_k}
 {\sum_j(\overline I_j-\underline I_j)},
 &\sum_j(\overline I_j-\underline I_j)>0,\\
 0,&\sum_j(\overline I_j-\underline I_j)=0.
 \end{cases}
\]

\section{PDE residual diagnostics}

Let
\[
 \mathcal R_\Phi(x)=\mathcal L[\Phi](x)-f(x)
\]
be the PDE residual function.  For every cell \(c\), let
\[
 [\mathcal R_\Phi]_c
 =[\underline{\mathcal R}_c,\overline{\mathcal R}_c]
\]
be a certified interval enclosure.  Define
\[
 M_c^{\mathcal R}=\magI([\mathcal R_\Phi]_c),\qquad
 w_c^{\mathcal R}=\wid([\mathcal R_\Phi]_c).
\]
The certified global pointwise residual upper bound is
\[
 \|\mathcal R_\Phi\|_{L^\infty}^{\rm upper}
 =\max_c M_c^{\mathcal R}.
\]
If available, also report
\[
 [\|\mathcal R_\Phi\|_{L^2}^2]
 \quad\text{and}\quad
 [\|\mathcal R_\Phi\|_{L^2}].
\]
Boundary-condition and initial-condition residuals must be represented by
separate residual functions and reported separately.

\section{Activation-enclosure diagnostics}

For cell \(c\), activation layer \(\ell\), neuron \(i\), and derivative order
\(r\in\{0,1,2\}\), let
\[
 [Z_{\ell i}]_c
 =[\underline Z_{c,\ell i},\overline Z_{c,\ell i}]
\]
be the certified preactivation interval used to construct the activation
enclosure.  Record
\[
 w^Z_{c,\ell i}=\wid([Z_{\ell i}]_c)
\]
and the certified approximation-error radius
\[
 \rho_{c,\ell i}^{(r)}.
\]
The associated approximation-error diameter is
\[
 d_{c,\ell i}^{(r)}=2\rho_{c,\ell i}^{(r)}.
\]

For comparisons within one layer and derivative order, define
\[
 S_\ell^{(r)}
 =\max_{c,i}\magI\left([\sigma^{(r)}(Z_{\ell i})]_c\right)
\]
and
\[
 \widehat\rho_{c,\ell i}^{(r)}
 =\begin{cases}
  \dfrac{\rho_{c,\ell i}^{(r)}}{S_\ell^{(r)}},
  &S_\ell^{(r)}>0,\\[2mm]
  0,&S_\ell^{(r)}=0.
 \end{cases}
\]
This is the canonical normalized activation approximation-error radius.

\subsection{Layerwise normalized interval-hull radius}

Let \([Q_{\ell i}]_c\) be the interval hull of a specified scalar neuron
quantity \(Q\) at layer \(\ell\).  The quantity identifier must be one of
\[
 Q\in\{Z,Y,\sigma^{(0)}(Z),\sigma^{(1)}(Z),\sigma^{(2)}(Z)\}.
\]
For fixed \(Q\) and \(\ell\), define
\[
 S_{Q,\ell}=\max_{c,i}\magI([Q_{\ell i}]_c)
\]
and
\[
 \nu_{c,\ell i}^{Q}
 =\begin{cases}
 \dfrac{\rad([Q_{\ell i}]_c)}{S_{Q,\ell}},&S_{Q,\ell}>0,\\[2mm]
 0,&S_{Q,\ell}=0.
 \end{cases}
\]
The medium benchmark must include the postactivation-value table
\(Q=Y\).  Tables for \(Z\), \(\sigma'(Z)\), and \(\sigma''(Z)\) are optional
extensions but must follow the same schema.

For a single wide table with layers as columns, aggregate over cells by the
volume-weighted mean
\[
 \bar\nu_{\ell i}^{Q}=\sum_c\omega_c\nu_{c,\ell i}^{Q}.
\]
The wide table entry in row \(i\), column \(\ell\) is
\(\bar\nu_{\ell i}^{Q}\).  If layer \(\ell\) has fewer than \(i\) neurons,
the entry is \(\NA\), not zero.

\section{Representation-complexity and reduction diagnostics}

For a sparse polynomial
\[
 P(\eps)=\sum_{\lambda\in\Lambda_P}P_\lambda\eps^\lambda,
\]
record
\[
 N_\alpha=p,\qquad N_\eta=q,\qquad
 N_{\rm mon}=|\Lambda_P|,
\]
\[
 d_{\max}=\max_{\lambda\in\Lambda_P}|\lambda|_1,
\]
\[
 d_{\alpha,\max}=\max_{\lambda\in\Lambda_P}|\lambda^\alpha|_1,
 \qquad
 d_{\eta,\max}=\max_{\lambda\in\Lambda_P}|\lambda^\eta|_1,
\]
and
\[
 N_{\alpha\eta}
 =\left|\left\{\lambda\in\Lambda_P:
 |\lambda^\alpha|_1>0\ \text{and}\ |\lambda^\eta|_1>0
 \right\}\right|.
\]
Duplicate exponent vectors must be merged before counting distinct monomials.
If \(P_\lambda\in U\) and \(\dim U=s\), the stored scalar coefficient count is
\[
 N_{\rm coeff}=s|\Lambda_P|.
\]

For support reduction with retained support \(\Lambda_{\rm keep}\) and
discarded support \(\Lambda_{\rm drop}\), record
\[
 N_{\rm before}=|\Lambda_P|,
\]
\[
 N_{\rm after}=|\Lambda_{\rm keep}|+N_{\rm new},
\]
\[
 R_{\rm support}
 =\begin{cases}
 1-N_{\rm after}/N_{\rm before},&N_{\rm before}>0,\\
 0,&N_{\rm before}=0,
 \end{cases}
\]
and, for a fixed coefficient norm,
\[
 M_{\rm drop}=\sum_{\lambda\in\Lambda_{\rm drop}}\|P_\lambda\|.
\]
When reduced and unreduced interval hulls are both available, define the
reduction-induced width increase
\[
 \Delta w_{\rm red}
 =\wid([P]_{\rm reduced})-\wid([P]_{\rm unreduced}).
\]

\section{Runtime, memory, and soundness diagnostics}

Record non-overlapping wall-clock times where applicable:
\[
 T_{\rm value},\ T_{\rm Jacobian},\ T_{\rm Hessian},\
 T_{\rm activation},\ T_{\rm multiplication},
\]
\[
 T_{\rm reduction},\ T_{\rm integration},\
 T_{\rm intervalization},\ T_{\rm total}.
\]
Also record peak allocated memory \(M_{\rm peak}\).  GPU measurements require
device synchronization before and after each timed region.

For sampled physical points \(x_n\in\mathcal X_c\), test containment of
function values and all derivatives that are part of the benchmark.  For a
scalar value \(v\) and interval \(I\), define
\[
 d(v,I)=\max\{\lb(I)-v,\ v-\ub(I),\ 0\}.
\]
Record the number of failures \(N_{\rm fail}\) and the maximum sampled
containment violation \(D_{\rm fail}\).  Sampling cannot prove soundness, but
any failure disproves the claimed enclosure.  Also record invalid-interval,
NaN-endpoint, and infinite-endpoint counts.

\section{Canonical data representation}

All benchmark data must be stored in UTF-8 comma-separated-value files with a
single header row, a period as decimal separator, and no thousands separator.
Missing values are represented by the literal string \(\NA\).  Boolean values
are represented by \texttt{0} and \texttt{1}.  Infinity is represented by
\texttt{inf} and negative infinity by \texttt{-inf}.  Each file must be
accompanied by a metadata file named \texttt{benchmark\_metadata.json}.

The metadata must contain at least:
\begin{verbatim}
{
  "schema_version": "1.2",
  "benchmark_level": "mini" or "medium",
  "problem_id": "...",
  "model_id": "...",
  "method_id": "...",
  "git_commit": "..." or null,
  "dtype": "float64",
  "device": "cpu" or "cuda:...",
  "quantile_interpolation": "linear",
  "cell_average": "volume_weighted" or "unweighted",
  "timestamp_utc": "ISO-8601 timestamp"
}
\end{verbatim}

\subsection{Canonical tidy metric table}

Every benchmark must contain \texttt{metrics.csv}.  Its exact column order is:
\begin{center}
\small
\begin{tabularx}{\textwidth}{@{}r l X@{}}
\toprule
Position & Column & Meaning \\
\midrule
1 & \texttt{schema\_version} & Schema version, currently \texttt{1.2}. \\
2 & \texttt{benchmark\_level} & \texttt{mini} or \texttt{medium}. \\
3 & \texttt{problem\_id} & PDE or benchmark identifier. \\
4 & \texttt{model\_id} & Trained-network identifier. \\
5 & \texttt{method\_id} & Certification-method identifier. \\
6 & \texttt{run\_id} & Unique identifier for the execution. \\
7 & \texttt{split\_id} & Domain partition or refinement identifier. \\
8 & \texttt{quantity} & Canonical quantity identifier. \\
9 & \texttt{metric} & Canonical metric identifier. \\
10 & \texttt{aggregation} & \texttt{none}, \texttt{mean\_weighted}, \texttt{max}, \texttt{q50}, \texttt{q90}, or \texttt{q99}. \\
11 & \texttt{derivative\_order} & \texttt{0}, \texttt{1}, \texttt{2}, or \NA. \\
12 & \texttt{layer} & Zero-based layer index or \NA. \\
13 & \texttt{neuron} & Zero-based neuron index or \NA. \\
14 & \texttt{output\_index} & Zero-based output component or \NA. \\
15 & \texttt{input\_index\_a} & First physical-input coordinate or \NA. \\
16 & \texttt{input\_index\_b} & Second coordinate for Hessians or \NA. \\
17 & \texttt{cell\_id} & Zero-based cell identifier or \NA. \\
18 & \texttt{value} & Numeric metric value. \\
19 & \texttt{unit} & Unit string, \texttt{dimensionless}, or \NA. \\
20 & \texttt{status} & \texttt{ok}, \texttt{missing}, \texttt{invalid}, or \texttt{not\_implemented}. \\
\bottomrule
\end{tabularx}
\end{center}

The pair \texttt{quantity}, \texttt{metric} determines the mathematical
meaning.  Canonical quantity identifiers include
\begin{center}
\texttt{Y}, \texttt{J}, \texttt{H}, \texttt{L2\_sq}, \texttt{L2},
\texttt{W12\_sq}, \texttt{W12}, \texttt{W22\_sq}, \texttt{W22},
\texttt{PDE\_residual}, \texttt{boundary\_residual},
\texttt{initial\_residual}, \texttt{complexity}, \texttt{runtime},
\texttt{memory}, and \texttt{soundness}.
\end{center}
Canonical metric identifiers include
\begin{center}
\texttt{lower}, \texttt{upper}, \texttt{width},
\texttt{local\_relative\_width}, \texttt{local\_relative\_radius},
\texttt{global\_normalized\_radius}, \texttt{frobenius\_width},
\texttt{relative\_norm\_width},
\texttt{domain\_volume\_normalized\_lower},
\texttt{domain\_volume\_normalized\_upper},
\texttt{domain\_volume\_normalized\_width},
\texttt{linf\_upper}, \texttt{n\_alpha}, \texttt{n\_eta},
\texttt{n\_monomials}, \texttt{n\_mixed\_monomials},
\texttt{max\_degree}, \texttt{seconds}, \texttt{bytes},
\texttt{failure\_count}, and \texttt{max\_violation}.
\end{center}

\subsection{Canonical cellwise interval table}

The mini and medium benchmarks must contain \texttt{cell\_intervals.csv} with
this exact column order:
\begin{center}
\small
\begin{tabularx}{\textwidth}{@{}r l X@{}}
\toprule
Position & Column & Meaning \\
\midrule
1 & \texttt{run\_id} & Execution identifier. \\
2 & \texttt{split\_id} & Partition identifier. \\
3 & \texttt{cell\_id} & Zero-based cell identifier. \\
4 & \texttt{cell\_weight} & \(\omega_c\), or \(1/N_{\rm cell}\) for unweighted output. \\
5 & \texttt{quantity} & \texttt{Y}, \texttt{J}, \texttt{H}, or residual identifier. \\
6 & \texttt{output\_index} & Output component or \NA. \\
7 & \texttt{input\_index\_a} & First derivative coordinate or \NA. \\
8 & \texttt{input\_index\_b} & Second derivative coordinate or \NA. \\
9 & \texttt{lower} & Interval lower endpoint. \\
10 & \texttt{upper} & Interval upper endpoint. \\
11 & \texttt{midpoint} & Interval midpoint. \\
12 & \texttt{radius} & Interval radius. \\
13 & \texttt{width} & Interval width. \\
14 & \texttt{magnitude} & Interval magnitude. \\
15 & \texttt{mignitude} & Interval mignitude. \\
16 & \texttt{local\_relative\_radius} & \(\rho_{\rm loc}\). \\
17 & \texttt{global\_normalized\_radius} & Familywise \(\rho_{\rm glob}\). \\
18 & \texttt{sign\_certified} & \texttt{1} if zero is excluded, otherwise \texttt{0}. \\
19 & \texttt{status} & \texttt{ok}, \texttt{invalid}, or \texttt{not\_implemented}. \\
20 & \texttt{local\_relative\_width} & \(\rho_{\rm wid}=2\rho_{\rm loc}\). \\
\bottomrule
\end{tabularx}
\end{center}
Rows are sorted lexicographically by
\[
 (\texttt{cell\_id},\texttt{quantity},\texttt{output\_index},
 \texttt{input\_index\_a},\texttt{input\_index\_b}).
\]

\subsection{Canonical norm table}

The benchmark must contain \texttt{norms.csv} with exact column order
\begin{center}
\small
\begin{tabularx}{\textwidth}{@{}r l X@{}}
\toprule
Position & Column & Meaning \\
\midrule
1 & \texttt{run\_id} & Execution identifier. \\
2 & \texttt{norm} & \texttt{L2}, \texttt{W12}, \texttt{W22}, or residual norm. \\
3 & \texttt{squared} & \texttt{1} for squared norm, otherwise \texttt{0}. \\
4 & \texttt{lower} & Certified lower endpoint. \\
5 & \texttt{upper} & Certified upper endpoint. \\
6 & \texttt{width} & Upper minus lower. \\
7 & \texttt{relative\_width} & Width divided by upper endpoint, with zero for \([0,0]\). \\
8 & \texttt{value\_contribution\_upper} & \(C_0^{\rm upper}\), or \NA. \\
9 & \texttt{gradient\_contribution\_upper} & \(C_1^{\rm upper}\), or \NA. \\
10 & \texttt{hessian\_contribution\_upper} & \(C_2^{\rm upper}\), or \NA. \\
11 & \texttt{value\_contribution\_width} & \(C_0^{\rm width}\), or \NA. \\
12 & \texttt{gradient\_contribution\_width} & \(C_1^{\rm width}\), or \NA. \\
13 & \texttt{hessian\_contribution\_width} & \(C_2^{\rm width}\), or \NA. \\
14 & \texttt{status} & \texttt{ok}, \texttt{invalid}, or \texttt{not\_implemented}. \\
15 & \texttt{domain\_volume} & Physical integration-domain volume \(V_{\mathcal X}\). \\
16 & \texttt{domain\_volume\_normalized\_lower} & Lower endpoint divided by \(V_{\mathcal X}\) for squared rows and by \(\sqrt{V_{\mathcal X}}\) otherwise. \\
17 & \texttt{domain\_volume\_normalized\_upper} & Upper endpoint with the same scaling. \\
18 & \texttt{domain\_volume\_normalized\_width} & Width with the same scaling. \\
\bottomrule
\end{tabularx}
\end{center}
Rows are ordered first by \texttt{norm} in the order
\texttt{L2}, \texttt{W12}, \texttt{W22}, followed by other norms, and then by
\texttt{squared} with squared rows first.

\section{Mini benchmark}

The mini benchmark is the mandatory smallest benchmark.  It is intended for
fast regression tests and standardized comparisons of certification methods.
It must be possible to execute it without storing all intermediate polynomial
objects.

\subsection{Required outputs}

The mini benchmark must produce:
\begin{enumerate}[label=\arabic*.]
 \item \texttt{benchmark\_metadata.json};
 \item \texttt{metrics.csv};
 \item \texttt{cell\_intervals.csv};
 \item \texttt{norms.csv};
 \item \texttt{complexity.csv};
 \item \texttt{timings.csv};
 \item \texttt{soundness.csv}.
\end{enumerate}

\subsection{Norm-output program contract}

The following contract is part of the definition of both benchmark programs.
It is not an optional reporting convention.

For every implemented norm
\[
 N\in\{\texttt{L2},\texttt{W12},\texttt{W22}\},
\]
the benchmark program must emit both the squared and unsquared certified norm
rows in \texttt{norms.csv}.  Every such row with status \texttt{ok} must
populate all of the following fields:
\[
 \texttt{lower},\quad \texttt{upper},\quad \texttt{width},\quad
 \texttt{relative\_width},\quad \texttt{domain\_volume},
\]
\[
 \texttt{domain\_volume\_normalized\_lower},\quad
 \texttt{domain\_volume\_normalized\_upper},\quad
 \texttt{domain\_volume\_normalized\_width}.
\]
For a squared norm row, the three domain-volume-normalized endpoints and width
are obtained by division by \(V_{\mathcal X}\).  For an unsquared norm row,
they are obtained by division by \(\sqrt{V_{\mathcal X}}\).

A program that reports only the unnormalized interval, only the relative width,
or only a scalar domain-volume-normalized point estimate does not implement the
mini benchmark.  Since the medium benchmark inherits the complete mini
benchmark, such a program also does not implement the medium benchmark.


\subsection{Required diagnostics}

For final function values \(Y\), report:
\begin{itemize}
 \item volume-weighted mean width;
 \item maximum width;
 \item \(0.50,0.90,0.99\) width quantiles;
 \item volume-weighted mean global normalized radius;
 \item maximum global normalized radius.
\end{itemize}

For the final Jacobian \(J\), report the same entrywise summaries and also
\[
 \Mean_\omega(W_F^J),\qquad \max_c W_{c,F}^J.
\]

For the final Hessian \(H\), report the corresponding diagnostics whenever
second derivatives are propagated or required by the PDE.  Otherwise rows must
be present in \texttt{metrics.csv} with status \texttt{not\_implemented}; they
must not be silently omitted.

For each implemented function-space norm, report the squared and unsquared
intervals, absolute widths, relative norm widths, the physical domain volume,
and the corresponding domain-volume-normalized intervals.  The minimum
standard set is \(L^2\) and \(W^{1,2}\); \(W^{2,2}\) is required when a
Hessian is propagated.

\paragraph{Mandatory norm-row contract.}
For every required norm name \(N\in\{\texttt{L2},\texttt{W12}\}\), and
also \(N=\texttt{W22}\) whenever the Hessian is propagated, the mini
benchmark must write exactly two rows to \texttt{norms.csv}: first the squared
row \(\texttt{squared}=1\), then the unsquared row
\(\texttt{squared}=0\).  In both rows, columns 15--18 are mandatory and
must contain
\[
 V_{\mathcal X},\qquad
 \frac{\underline{\mathcal N^2}}{V_{\mathcal X}},\quad
 \frac{\overline{\mathcal N^2}}{V_{\mathcal X}},\quad
 \frac{\overline{\mathcal N^2}-\underline{\mathcal N^2}}
      {V_{\mathcal X}}
\]
for the squared row, and
\[
 V_{\mathcal X},\qquad
 \frac{\underline{\mathcal N}}{\sqrt{V_{\mathcal X}}},\quad
 \frac{\overline{\mathcal N}}{\sqrt{V_{\mathcal X}}},\quad
 \frac{\overline{\mathcal N}-\underline{\mathcal N}}
      {\sqrt{V_{\mathcal X}}}
\]
for the unsquared row.  These values must not be omitted, replaced by
\texttt{NA}, or inferred only by downstream plotting code when the underlying
norm interval has status \texttt{ok}.  A mini-benchmark implementation is
non-conforming if these columns are absent or unpopulated.

For PDE certification, report
\[
 \|\mathcal R_\Phi\|_{L^\infty}^{\rm upper}
\]
and, when implemented, certified \(L^2\) residual intervals.  Boundary and
initial residuals are separate quantities.

For representation complexity, record at least
\[
 N_\alpha,\quad N_\eta,\quad N_{\rm mon},\quad
 N_{\alpha\eta},\quad d_{\max},\quad N_{\rm coeff}
\]
for \(Y\), \(J\), and \(H\) separately.

For computational cost, record total runtime, the principal non-overlapping
stage times, and peak memory.  For soundness diagnostics, record
\(N_{\rm fail}\), \(D_{\rm fail}\), invalid interval count, NaN endpoint count,
and infinite endpoint count.

\subsection{Canonical mini-benchmark summary order}

For human-readable printing, the summary rows must appear in this order:
\begin{enumerate}[label=\arabic*.]
 \item run metadata and domain partition;
 \item function-value enclosure diagnostics;
 \item Jacobian enclosure diagnostics;
 \item Hessian enclosure diagnostics;
 \item squared and unsquared function-space norm intervals, immediately
 followed by their domain-volume-normalized lower endpoint, upper endpoint,
 and width;
 \item PDE, boundary, and initial residual diagnostics;
 \item representation-complexity diagnostics for \(Y,J,H\);
 \item runtime and peak-memory diagnostics;
 \item soundness and validity diagnostics.
\end{enumerate}
Within each enclosure family, print metrics in the order
\[
 \texttt{mean\_width},\ \texttt{max\_width},\
 \texttt{q50\_width},\ \texttt{q90\_width},\ \texttt{q99\_width},
\]
\[
 \texttt{mean\_global\_normalized\_radius},\
 \texttt{max\_global\_normalized\_radius},
\]
followed by family-specific matrix or tensor diagnostics.

\section{Medium benchmark}

The medium benchmark contains every mini-benchmark output, file, row, column,
and diagnostic without modification.  In particular, it must produce the same
\texttt{norms.csv} required by the mini benchmark, including the mandatory
squared and unsquared domain-volume-normalized norm intervals in columns
15--18.  The medium benchmark is therefore not permitted to omit, rename, or
replace these quantities with only relative widths.  It additionally diagnoses
where approximation error enters the network and how interval hulls evolve
through the hidden layers.

\paragraph{Inheritance rule.}
A medium-benchmark program must first satisfy every mini-benchmark conformance
condition.  Its additional per-neuron and layerwise files extend the mini
benchmark; they do not form an alternative output schema.  Consequently, a
medium-benchmark run is non-conforming whenever the corresponding mini
benchmark would be non-conforming, including whenever any required
\texttt{domain\_volume\_normalized\_lower},
\texttt{domain\_volume\_normalized\_upper}, or
\texttt{domain\_volume\_normalized\_width} entry is missing from an
implemented norm row.

\subsection{Per-neuron activation approximation table}

The medium benchmark must contain
\texttt{activation\_approximation.csv}.  Its exact column order is:
\begin{center}
\small
\begin{tabularx}{\textwidth}{@{}r l X@{}}
\toprule
Position & Column & Meaning \\
\midrule
1 & \texttt{run\_id} & Execution identifier. \\
2 & \texttt{split\_id} & Partition identifier. \\
3 & \texttt{cell\_id} & Zero-based cell identifier. \\
4 & \texttt{cell\_weight} & Physical-volume weight \(\omega_c\). \\
5 & \texttt{layer} & Zero-based hidden activation-layer index. \\
6 & \texttt{neuron} & Zero-based neuron index within the layer. \\
7 & \texttt{derivative\_order} & \texttt{0}, \texttt{1}, or \texttt{2}. \\
8 & \texttt{preactivation\_lower} & \(\underline Z_{c,\ell i}\). \\
9 & \texttt{preactivation\_upper} & \(\overline Z_{c,\ell i}\). \\
10 & \texttt{preactivation\_midpoint} & Midpoint of \([Z_{\ell i}]_c\). \\
11 & \texttt{preactivation\_radius} & Radius of \([Z_{\ell i}]_c\). \\
12 & \texttt{preactivation\_width} & Width of \([Z_{\ell i}]_c\). \\
13 & \texttt{approximation\_kind} & Canonical approximation identifier, e.g. \texttt{affine}, \texttt{quadratic}, or \texttt{interval}. \\
14 & \texttt{approximation\_error\_radius} & \(\rho_{c,\ell i}^{(r)}\). \\
15 & \texttt{approximation\_error\_diameter} & \(2\rho_{c,\ell i}^{(r)}\). \\
16 & \texttt{activation\_scale} & \(S_\ell^{(r)}\). \\
17 & \texttt{normalized\_approximation\_radius} & \(\widehat\rho_{c,\ell i}^{(r)}\). \\
18 & \texttt{noise\_symbol\_id} & Stable identifier for \(\eta_{\ell i}^{(r)}\), or \NA if no fresh symbol is introduced. \\
19 & \texttt{shared\_noise\_group} & Identifier when one approximation noise symbol is reused, otherwise \NA. \\
20 & \texttt{status} & \texttt{ok}, \texttt{invalid}, or \texttt{not\_implemented}. \\
\bottomrule
\end{tabularx}
\end{center}
Rows are sorted by
\[
 (\texttt{cell\_id},\texttt{layer},\texttt{neuron},
 \texttt{derivative\_order}).
\]
The derivative-order rows \(0,1,2\) must be present whenever the corresponding
activation enclosure is used by the propagated value, Jacobian, or Hessian.

\subsection{Layerwise normalized interval-radius table}

The medium benchmark must contain a separate wide table named
\begin{center}\path{layer_normalized_radius_Y.csv}.\end{center}
It summarizes the postactivation
interval hulls \([Y_{\ell i}]_c\).  Its exact column order is
\[
 \texttt{neuron},\ \texttt{layer\_0},\ \texttt{layer\_1},\ \ldots,
 \texttt{layer\_{L-1}}.
\]
The row index \texttt{neuron} is zero-based.  The entry in row \(i\), column
\texttt{layer\_\(\ell\)} is
\[
 \bar\nu_{\ell i}^{Y}
 =\sum_c\omega_c
 \begin{cases}
 \dfrac{\rad([Y_{\ell i}]_c)}
 {\max_{c',j}\magI([Y_{\ell j}]_{c'})},
 &\max_{c',j}\magI([Y_{\ell j}]_{c'})>0,\\[3mm]
 0,&\text{otherwise}.
 \end{cases}
\]
If layer \(\ell\) has fewer than \(i+1\) neurons, the entry is \(\NA\).
Columns always follow increasing layer index; rows always follow increasing
neuron index.

The medium benchmark may additionally provide the same wide schema for other
quantities using the filenames
\begin{center}
\texttt{layer\_normalized\_radius\_Z.csv},\\
\texttt{layer\_normalized\_radius\_sigma1.csv},\\
\texttt{layer\_normalized\_radius\_sigma2.csv}.
\end{center}
Here \texttt{sigma1} and \texttt{sigma2} refer to interval hulls of
\(\sigma'(Z_{\ell i})\) and \(\sigma''(Z_{\ell i})\), respectively.  These
optional tables must use the identical row and column conventions.

\subsection{Required medium-benchmark visual summaries}

A medium-benchmark report should render, from the canonical CSV data:
\begin{enumerate}[label=\arabic*.]
 \item one table per derivative order \(r=0,1,2\), with rows corresponding to
 neurons and grouped columns containing preactivation interval, approximation
 kind, approximation-error radius, and normalized approximation radius;
 \item the wide layer table \texttt{layer\_normalized\_radius\_Y.csv};
 \item a ranked list of the largest approximation-error radii;
 \item a ranked list of the largest normalized approximation-error radii;
 \item the worst cells and neurons according to preactivation width and
 approximation-error radius.
\end{enumerate}
These visual summaries are derived views.  The CSV files remain the canonical
data representation.

\section{Implementation invariants and safety checks}

An implementation conforming to this reference must satisfy the following
invariants.

\begin{enumerate}[label=\arabic*.]
 \item Every reported metric is computable from explicitly stored endpoints,
 coefficients, counts, or timings.  No hidden stabilization constant is used.
 \item The only polynomial variable classes are domain noise symbols \(\alpha\)
 and approximation noise symbols \(\eta\).
 \item The word \emph{residual} denotes an actual function difference, never a
 noise symbol.
 \item All derivative indices refer to derivatives with respect to the physical
 input \(x\).
 \item Missing or unimplemented quantities are represented by \(\NA\) and an
 explicit non-\texttt{ok} status; they are never silently replaced by zero.
 \item Zero is used only when the mathematical definition prescribes zero, such
 as the normalized radius of an identically zero interval family.
 \item Cell-weighted means use the stored \texttt{cell\_weight} values, which
 must sum to one within numerical tolerance.
 \item Distinct monomials are counted only after duplicate exponent vectors are
 merged.
 \item Squared norm intervals are retained even when unsquared norm intervals
 are also reported.
 \item Sampling diagnostics are labelled as diagnostics and are not described
 as proofs of soundness.
\end{enumerate}

\section{Canonical vocabulary}

\renewcommand{\arraystretch}{1.2}
\small
\begin{tabularx}{\textwidth}{@{}l X@{}}
\toprule
Preferred term & Meaning \\
\midrule
Interval hull & Final interval enclosure obtained from a polynomial enclosure over all domain and approximation noise symbols. \\
Absolute width & Upper endpoint minus lower endpoint. \\
Absolute radius & Half the absolute width. \\
Magnitude & Maximum absolute value of the interval endpoints. \\
Mignitude & Largest certified lower bound on absolute value over the interval. \\
Local relative width & Interval width divided by the magnitude of the same interval; equal to twice the local relative radius. \\
Local relative radius & Interval radius divided by the magnitude of the same interval. \\
Global normalized radius & Interval radius divided by the maximum magnitude over the same explicitly specified family. \\
Volume-weighted mean width & Componentwise mean width weighted by physical cell volume. \\
Relative norm width & Width of a nonnegative norm interval divided by its upper endpoint; this is a tightness metric. \\
Domain-volume-normalized norm & Norm divided by the square root of the physical integration-domain volume; this is a scale metric, not a tightness metric. \\
Approximation-error radius & Certified nonnegative scalar \(\rho\) in an activation enclosure. \\
Approximation noise symbol & A component of \(\eta\); this is the canonical name. \\
Approximation noise term & A product \(\rho\eta\). \\
Residual function & An actual difference between a function and its approximation, or between the two sides of a PDE. \\
Distinct monomial count & Cardinality of the canonical exponent support after duplicate exponents are merged. \\
Mixed monomial & A monomial depending on at least one domain noise symbol and at least one approximation noise symbol. \\
Reduction-induced width increase & Difference between reduced and unreduced interval-hull widths. \\
Mini benchmark & Mandatory compact diagnostic suite for routine regression and method comparison. \\
Medium benchmark & Complete mini benchmark plus per-neuron approximation diagnostics and layerwise interval-hull-radius tables. \\
\bottomrule
\end{tabularx}
\normalsize

\section{Default interpretation}

No normalized metric replaces the absolute width.  The default pair for
pointwise enclosures is
\[
 \boxed{\text{absolute width}}
 \qquad\text{and}\qquad
 \boxed{\text{global normalized radius}}.
\]
The local relative radius and the equivalent local relative width are secondary
diagnostics.  For nonnegative function-space norm intervals, the default
tightness metric is the relative norm width
\[
 \rho_{\mathcal N}
 =\frac{\overline{\mathcal N}-\underline{\mathcal N}}
 {\overline{\mathcal N}}
\]
when \(\overline{\mathcal N}>0\), and zero for the exact interval \([0,0]\).
The domain-volume-normalized norm interval is reported separately whenever
\(V_{\mathcal X}\) is available.  It describes average physical scale over
the domain and must not be interpreted as enclosure tightness.

The canonical CSV schemas, fixed column orders, benchmark-level requirements,
and implementation invariants are part of the mathematical reference.  Agents
must not change them silently.  Any extension must preserve the existing
columns and identifiers and must increase the schema version when semantics
change.

\end{document}
